% capitolo 2: descrizione generale dello stage
\chapter{Progetto di stage}
\label{cap:progetto-di-stage}

\intro{In questo capitolo viene introdotto il progetto di stage, descrivendone le motivazioni iniziali, gli obiettivi prefissati, i vincoli e la pianificazione delle fasi di lavoro.}\\

\section{L'idea}

\subsection{Motivazioni: perchè adottare un Data Catalog}
Nell'era moderna, le aziende hanno un'enormità di dati sparsi in sistemi diversi (database, fogli Excel, file CSV...). Senza un catalogo, le persone 
non sanno cosa esiste, dove si trova o se possono fidarsi di quel dato. Il Data Catalog risolve questo caos creando una mappa unica, organizzata e facilmente consultabile.

\subsubsection{Sicurezza e compliance normativa}
Storicamente l'individuazione dei dati personali ai fini GDPR avveniva in modo manuale, producendo un semplice "catalogo grezzo". 
L'adozione di un Data Catalog moderno automatizza questo processo scansionando i database e classificando i dati sensibili in autonomia.
Il passaggio da un approccio manuale a un catalogo completo e automatizzato segna una forte crescita della maturità aziendale nella sicurezza e nel modello \gls{samm}. 
Grazie a questa mappatura esatta e dinamica, l'azienda sa sempre dove e come applicare le necessarie policy di protezione, come il mascheramento o la pseudonimizzazione.

\subsubsection{Analisi dei dati e Business Intelligence}
Gli analisti accedono al catalogo per cercare le informazioni di cui hanno bisogno per creare i loro report. 
Se un analista ha bisogno del fatturato del 2023, cerca "Fatturato" nel Data Catalog: il sistema gli dice subito in quale 
database si trova, se il dato è aggiornato e di alta qualità, eliminando il rischio di errori o di ricreare set di dati che un collega aveva già preparato in passato.

\subsubsection{Data Quality e Governance}
Nel catalogo è possibile valutare la qualità delle informazioni, vedere chi è il responsabile di una specifica tabella e capire come i dati sono collegati tra loro. 
Assicura che chiunque in azienda prenda decisioni basandosi su informazioni affidabili, certificate e sicure, migliorando l'efficienza generale.


\section{Obiettivi dello stage}
% obiettivi dell'azienda (vedi prog. formativo) + altre cose tipo: studio completo del way of working aziendale per un mio possibile futuro in azienda
% studio delle tecnologie, dell'ambiente, confidenza con strumenti, competenze trasversali in più ambiti per tenere tutte le porte aperte

\subsection{Obiettivi personali}
% tipo aspettative personali e cosa mi aspetto/voglio dallo stage


\section{Vincoli di progetto}
\subsection{Vincoli tecnologici}
% per esempio le tecnologie che ho usato e che ho dovuto installare

\subsection{Vincoli metodologici}
% usare Teams per qualsiasi comubnicazione, studio della comunicazione (vedi doc), metodologia agile, aggironamenti frequenti, 
% scrittura di tutta la documentazione richiesta: dai documenti ufficiali per lo studio di OpenMetadata, ai verbali delle riunioni, a quelli relativi alla formazione, come
% i manuali di installazione o i verbali sugli incontri su Power BI

\subsection{Vincoli temporali}
% le ore a disposizione, considerando che il venerdì si esce un'ora prima e rispettando le pause. I vincoli temporali poi sono stati definiti nel piano di progetto che riporto dopo

\section{Pianificazione del lavoro}


\section{Formazione e tecnologie}
% formazione su SQL e database: query, sintassi SQL base, viste anche stored procedures, viste, cursori; confidenza con PostgreSQL e PgAdmin e anche con SQL Server; procedure di backup e restore
% docker e docker compose, WSL, VM (Virtual Box), perchè il pc è windows ma mi serviva anche ubuntu (+ studio comandi batch ubuntu)
% formazione su sviluppo sicuro, principi della privacy e protezione dei dati sensibili e personali
% altro: metodologia Agile (SCRUM), git e github, sistema gestionale di Wintech (ESOLVER), Power BI e principi della business intelligence




